\section{Guessing the active set instead of walking to it}
\label{sec:pdas}

The dual method reaches the active set one constraint at a time. That is its
defining feature and, at moderate $n$, its principal cost: a budget-plus-bounds
problem in $100$ variables takes some $74$ outer iterations, each of which is a
handful of matrix--vector products. The alternative is to guess the whole active
set at once, solve the working-set subproblem it induces, and repair the guess
from the signs that come back. This is the semismooth Newton method of
Hinterm\"uller, Ito and Kunisch~\cite{hintermuller2002} read as a primal--dual
active-set strategy, and block principal pivoting on the associated linear
complementarity problem~\cite{judice1994,murty1988,cottle1992} read as a rule for
exchanging several indices at a time.

\subsection{The iteration}

Given a candidate set $\Active$, Proposition~\ref{prop:sub} gives the working-set
solution directly:
\begin{equation}
  \bigl(C_\Active\T G^{-1} C_\Active\bigr)\lambda_\Active
   = b_\Active - C_\Active\T x_u,
  \qquad
  x = x_u + G^{-1}C_\Active \lambda_\Active,
  \qquad x_u = G^{-1}a,
  \label{eq:pdassolve}
\end{equation}
with $\lambda_i = 0$ for $i \notin \Active$. Both solves go through one Cholesky
factorisation of $G$, computed once, and one of the $k \times k$ dual Hessian,
computed per repair. The repair rule reads the two sign conditions of
\eqref{eq:dfeas}--\eqref{eq:pfeas} as instructions:
\begin{equation}
  \Active' = \{1,\dots,m_{\mathrm{eq}}\} \;\cup\;
  \bigl\{\, i > m_{\mathrm{eq}} \;:\;
    (i \in \Active \text{ and } \lambda_i \ge 0)
    \ \text{ or }\
    (c_i\T x < b_i) \,\bigr\},
  \label{eq:repair}
\end{equation}
that is: an active constraint whose multiplier has gone negative does not belong
in the set and leaves it; an inactive constraint that is violated at $x$ does
belong and joins it; equalities are always in. The starting guess is the
equalities together with whatever the unconstrained minimiser violates, which is
already the answer when it violates nothing. If $\Active' = \Active$ the
conditions~\eqref{eq:pfeas}--\eqref{eq:comp} hold to the tolerance used, and
\eqref{eq:pdassolve} supplies \eqref{eq:stat} by construction.

The whole set is exchanged at once, which is what converges in two to four repairs
independently of $n$. It trades arithmetic, which is abundant at these sizes, for
iterations, which are not.

\subsection{Why there is no finite-termination theorem here}

Block exchange can cycle, and the standard remedy is a least-index rule --- flip
only the lowest-indexed offending constraint, in the manner of Bland's rule for
the simplex method and Murty's least-index rule for linear complementarity
problems~\cite{murty1988}. In the bound-constrained case that remedy comes with a
theorem, and it is worth being precise about why the theorem does not survive the
generalisation, because the reason is structural rather than a gap in the
argument.

Consider $\min_{x \ge 0} \tfrac12 x\T G x - a\T x$, so that $C = I$ and $m = n$.
The system solved on a candidate set $\Active$ is then
\begin{equation}
  \bigl(I_{\cdot\Active}\T G^{-1} I_{\cdot\Active}\bigr)\lambda_\Active
  = \bigl(G^{-1}\bigr)_{\Active\Active}\lambda_\Active
  = b_\Active - x_{u,\Active},
  \label{eq:principal}
\end{equation}
whose coefficient matrix is a \emph{principal submatrix} of $G^{-1}$. Since
$G^{-1} \succ 0$, every principal submatrix of it is positive definite, so
$G^{-1}$ is a $P$-matrix: every principal pivot is well defined, whatever the
guess. That is the hypothesis under which block principal pivoting with a
least-index guard terminates finitely at the unique
solution~\cite{judice1994,murty1988}, and the guarantee survives inexact inner
solves once their residual is small against the problem's decision
margin~\cite{schmelzer2026nncg}.

None of it is available for general $C$. The matrix
$C_\Active\T G^{-1} C_\Active$ is positive definite exactly when $C_\Active$ has
full column rank, and that is a property of the guess, not of the data: a guessed
set may name $k > n$ constraints, or $k \le n$ linearly dependent ones, and then
the pivot is undefined rather than merely awkward. There is no $P$-matrix to
appeal to, and correspondingly no theorem to inherit. Contrast
Proposition~\ref{prop:rank}: the dual method's step rules make dependence
unreachable, and it is exactly that protection which guessing forfeits.

The implementation's response is to treat the Cholesky factorisation of
$C_\Active\T G^{-1}C_\Active$ as the rank test --- a dependent guess fails there
rather than returning something plausible --- and to treat the least-index rule as
a way of making progress where block exchange oscillates rather than as a
termination proof. Neither is a substitute for the guarantee. The certificate is.

\subsection{The certificate}

Because the method may fail, and may fail by stabilising on a set that is simply
not optimal, the trade is sound only because the answer is checked. This is where
Proposition~\ref{prop:sufficient} earns its place: for a strictly convex program
the KKT conditions are sufficient, so a candidate $(x,\lambda)$ satisfying
\begin{equation}
  \norm{G x - a - C\lambda}_\infty \le \varepsilon,
  \quad
  |c_i\T x - b_i| \le \varepsilon \ (i \le m_{\mathrm{eq}}),
  \quad
  c_i\T x - b_i \ge -\varepsilon, \ \ \lambda_i \ge -\varepsilon,
  \ \ |\lambda_i (c_i\T x - b_i)| \le \varepsilon
  \label{eq:certificate}
\end{equation}
for the inequalities, with $\varepsilon$ scaled by the data, \emph{is} the answer
up to that tolerance --- not a plausible candidate awaiting a second opinion.
Stationarity is checked against $G$ directly rather than trusted from the
construction, since the construction is precisely what an ill-conditioned working
set corrupts. A candidate failing~\eqref{eq:certificate} is discarded and the
exact walk of Algorithm~\ref{alg:gi} is run instead, so the guess can never
degrade the answer: it either produces the minimiser or produces nothing.

\begin{remark}[Two tolerances with very different jobs]
The tolerance in~\eqref{eq:repair} decides which set is tried next. A bad choice
there costs a repair or a fallback and can never cost correctness, so it is not
delicate. The tolerance in~\eqref{eq:certificate} is the only thing standing
between a non-optimal point and the caller, and is therefore the one to reason
about. The gap it must separate is empirically wide: over the implementation's
measurements, converged attempts that were optimal satisfied
\eqref{eq:certificate} with residual at most $4\times10^{-15}$, while those that
were not missed by $10^{-1}$ or worse. Any threshold in between separates them,
which is the sense in which the certificate is robust rather than tuned.
\end{remark}
